\documentclass[11pt, a4paper]{article}

\usepackage[T1]{fontenc}
\usepackage[utf8]{inputenc}
\usepackage{tgpagella}
\usepackage[spanish, es-noshorthands]{babel}
\usepackage{amsmath, amssymb, bm}
\usepackage{booktabs}
\usepackage{tabularx}
\usepackage[table, xcdraw]{xcolor}
\usepackage[most]{tcolorbox}
\usepackage[american]{circuitikz}
\usepackage[margin=2.5cm]{geometry}
\usepackage{fancyhdr}

\pagestyle{fancy}
\fancyhf{}
\setlength{\headheight}{16pt}
\lhead{\textbf{UCB Sede Tarija} | Ing. Mecatrónica}
\chead{}
\rhead{IMT-121 Circuitos Electrónicos I | Ejemplo Guiado S06-3}
\lfoot{Prof: M.Sc. Ing. Bernardo Quiroga Turdera}
\rfoot{Página \thepage}
\renewcommand{\headrulewidth}{0.4pt}

\definecolor{ucbblue}{RGB}{0, 51, 102}

\begin{document}

\begin{tcolorbox}[colback=ucbblue!5!white, colframe=ucbblue, title=\textbf{Hoja de Trabajo Guiada: Thévenin y Norton}]
    \centering \large \textbf{Reducción de Redes desde un Puerto de Dos Terminales}
\end{tcolorbox}

\section*{Circuito de Estudio}
El objetivo es simplificar la "Red Interna" de la izquierda y encontrar sus circuitos equivalentes de Thévenin y Norton respecto a los terminales del puerto $a-b$, donde eventualmente conectaremos la carga $R_L$.

\begin{center}
\begin{circuitikz}[scale=0.85, transform shape, thick]
        % Red interna
        \draw (0,0) to[V, l={$12\,\text{V}$}] (0,2) to[R, l={$4\,\text{k}\Omega$}] (3,2);
        \draw (2.5,2) to[short, *-o] (6,2) node[above]{$a$};
        \draw (2.5,2) to[R, l_={$8\,\text{k}\Omega$}] (2.5,0);
        \draw (4.5,0) to[I, l={$1\,\text{mA}$}] (4.5,2);
        
        \draw (0,0) to[short, -o] (6,0) node[below]{$b$};
        \draw (2.5,0) node[ground]{};
        
        % Carga (removible)
        \draw[dashed, blue] (5,2) -- (6,2) to[R, l={$R_L$}] (6,0) -- (5,0);
    \end{circuitikz}
\end{center}

\section*{Paso 1: Aislar el Puerto y Hallar $V_{th}$}
\textbf{Definición:} $V_{th}$ es el voltaje en los terminales $a-b$ cuando están en circuito abierto.
\begin{enumerate}
    \item Retira (borra) mentalmente la carga $R_L$.
    \item Calcula la tensión en el nodo $a$. (Tip: Puedes usar superposición para simplificar el análisis).
\end{enumerate}
\vspace{3cm}

\section*{Paso 2: Desactivar Fuentes y Hallar $R_{th}$}
\textbf{Definición:} $R_{th}$ es la resistencia que se mide "mirando" hacia adentro desde el puerto $a-b$ cuando todas las fuentes independientes han sido apagadas.
\begin{enumerate}
    \item Redibuja el circuito: Las fuentes de tensión ideales se reemplazan por un \underline{\hspace{3cm}}, y las fuentes de corriente por un \underline{\hspace{3cm}}.
    \item Dibuja una flecha entrando por $a-b$ y calcula la resistencia equivalente.
\end{enumerate}
\vspace{4.5cm}

\section*{Paso 3: Calcular Norton y Dibujar los Equivalentes}
Sabiendo que $R_N = R_{th}$ y que $I_N = \frac{V_{th}}{R_{th}}$, calcula el modelo de Norton. \\
Luego, dibuja ambos modelos equivalentes \textbf{conectados a la carga $R_L = 8\,\text{k}\Omega$}.

\vspace{5cm}
\begin{itemize}
    \item \textbf{Verificación Final:} Para el equivalente de Thévenin, ¿cuánto vale $I_L$? ¿Cuánto vale $V_L$?
\end{itemize}

\end{document}
